\documentclass[conference]{IEEEtran}
\IEEEoverridecommandlockouts

\usepackage{cite}
\usepackage{amsmath,amssymb,amsfonts}
\usepackage{graphicx}
\usepackage{url}

\def\BibTeX{{\rm B\kern-.05em{\sc i\kern-.025em b}\kern-.08em
    T\kern-.1667em\lower.7ex\hbox{E}\kern-.125emX}}

\begin{document}

\title{Machine Learning Classification of Heart and Lung Sound Recordings}

\author{
\IEEEauthorblockN{Miguel Diaz-Alvarez, Jordan Mckoy, and AJ Carroll}
\IEEEauthorblockA{
\textit{Department of Electrical and Computer Engineering}\\
\textit{University of North Carolina at Charlotte}\\
Charlotte, NC, USA}
}

\maketitle

\begin{abstract}
This report presents a machine learning approach for classifying heart, lung,
and mixed audio recordings. The project uses the HLS-CMDS dataset, which
contains separate heart sound recordings, lung sound recordings, and mixed
heart/lung recordings. The workflow includes audio preprocessing, feature
extraction, model training, hyperparameter tuning, and performance evaluation.
The final model is evaluated using classification metrics such as accuracy,
precision, recall, F1-score, and confusion matrices. Replace this paragraph
with a concise summary of the completed method, best result, and main finding.
\end{abstract}

\begin{IEEEkeywords}
machine learning, audio classification, heart sounds, lung sounds, feature
extraction, neural networks
\end{IEEEkeywords}

\section{Introduction}
The goal of this project is to classify biomedical audio recordings into
meaningful sound categories using machine learning. Heart and lung sounds are
clinically important signals, and automated classification can support faster
screening, data organization, and diagnostic decision support.

Briefly introduce the problem, why it matters, and what this project attempts
to solve. End the section with a short summary of the main contributions, such
as:

\begin{itemize}
    \item preprocessing and organizing the HLS-CMDS audio dataset;
    \item extracting useful audio features for classification;
    \item training and comparing machine learning models;
    \item evaluating model performance using standard classification metrics.
\end{itemize}

\section{Dataset}
The dataset used in this project is the HLS-CMDS dataset. It contains heart
sound recordings, lung sound recordings, and mixed heart/lung recordings. The
project directory is organized into the HS, LS, and Mix subsets.

\begin{table}[htbp]
\caption{Dataset Summary}
\begin{center}
\begin{tabular}{lcc}
\hline
\textbf{Subset} & \textbf{Description} & \textbf{Recordings} \\
\hline
HS & Heart sound recordings & 50 \\
LS & Lung sound recordings & 50 \\
Mix & Mixed heart/lung recordings & 435 \\
\hline
\end{tabular}
\label{tab:dataset-summary}
\end{center}
\end{table}

Describe the source of the data, the class labels, the number of samples per
class, and any relevant characteristics of the audio files, such as sampling
rate, duration, or class imbalance.

\section{Preprocessing and Feature Extraction}
Before model training, the audio recordings were preprocessed to produce a
consistent input representation. This section should describe the steps used to
clean, normalize, transform, or segment the audio data.

Possible preprocessing steps include:

\begin{itemize}
    \item loading waveform files and labels;
    \item resampling audio to a consistent sampling rate;
    \item trimming or padding recordings to a fixed length;
    \item normalizing signal amplitude;
    \item converting audio into features such as MFCCs, spectrograms, or
    Mel-spectrograms.
\end{itemize}

If feature extraction was used, include the feature dimensions and explain why
those features are appropriate for this classification task.

\section{Methodology}
This section describes the machine learning models and training procedure used
for the project. Include enough detail for another person to reproduce the
experiment.

\subsection{Model Architecture}
Describe the model or models tested. For example, if a convolutional neural
network was used, describe the convolutional layers, activation functions,
pooling layers, fully connected layers, dropout, and output layer.

\subsection{Training Procedure}
Describe the train-validation-test split, loss function, optimizer, learning
rate, batch size, number of epochs, and any hyperparameter tuning strategy.

\subsection{Evaluation Metrics}
The trained models were evaluated using standard classification metrics:

\begin{itemize}
    \item accuracy;
    \item precision;
    \item recall;
    \item F1-score;
    \item confusion matrix.
\end{itemize}

\section{Results}
Present the main experimental results in this section. Include a table comparing
the tested models or configurations.

\begin{table}[htbp]
\caption{Model Performance Comparison}
\begin{center}
\begin{tabular}{lcccc}
\hline
\textbf{Model} & \textbf{Accuracy} & \textbf{Precision} &
\textbf{Recall} & \textbf{F1-score} \\
\hline
Baseline model & TODO & TODO & TODO & TODO \\
Tuned model & TODO & TODO & TODO & TODO \\
Final model & TODO & TODO & TODO & TODO \\
\hline
\end{tabular}
\label{tab:model-results}
\end{center}
\end{table}

Add the most important figures here, such as a confusion matrix, training loss
curve, validation accuracy curve, or feature visualization.

\begin{figure}[htbp]
\centerline{%
\IfFileExists{figures/confusion_matrix.png}
{\includegraphics[width=0.45\textwidth]{figures/confusion_matrix.png}}
{\fbox{\parbox[c][1.25in][c]{0.42\textwidth}{\centering Upload
\texttt{figures/confusion\_matrix.png} or replace this placeholder.}}}}
\caption{Confusion matrix for the final classification model. Replace the image
path with the exported figure from the project.}
\label{fig:confusion-matrix}
\end{figure}

\section{Discussion}
Discuss what the results mean. Explain which model performed best, which
classes were easiest or hardest to classify, and whether the results suggest
overfitting, underfitting, or good generalization.

Also describe limitations of the project, such as dataset size, class imbalance,
audio quality, limited hyperparameter search, or constraints in preprocessing.

\section{Conclusion}
This project developed a machine learning pipeline for classifying heart, lung,
and mixed audio recordings. Summarize the final model, the best performance
achieved, and the main lesson learned from the project. End with possible future
work, such as collecting more data, improving feature extraction, testing deeper
models, or deploying the model in an application.

\section*{Acknowledgment}
Optional: thank the instructor, course staff, dataset providers, or teammates.
Delete this section if it is not needed.

\begin{thebibliography}{00}
\bibitem{b1} I. Goodfellow, Y. Bengio, and A. Courville, \textit{Deep Learning}.
Cambridge, MA, USA: MIT Press, 2016.

\bibitem{b2} B. McFee et al., ``librosa: Audio and music signal analysis in
Python,'' in \textit{Proc. 14th Python in Science Conf.}, 2015, pp. 18--25.

\bibitem{b3} F. Chollet et al., ``Keras,'' 2015. [Online]. Available:
\url{https://keras.io}

\bibitem{b4} A. Paszke et al., ``PyTorch: An imperative style, high-performance
deep learning library,'' in \textit{Advances in Neural Information Processing
Systems}, vol. 32, 2019.
\end{thebibliography}

\end{document}
